% GENERATED FILE. Edit canonical lessons, then run npm run generate:books.
% canonical-source-hash: fnv1a64:a66202562fab26ac
% canonical-lessons: ML-C436-laborattari, ML-C436-pariksanam, ML-C436-sarve, ML-C436-kantupitittam, ML-C436-sastriyamaya, ML-R436-review-429-436

\chapter{In the Lab}
\label{ch:ml-in-the-lab}
\hlchaptermodality{436}

\begin{chapteropening}
\textbf{By the end of this chapter:} \emph{I can say a laboratory, an experiment, a survey, a discovery and scientific.}

Complete the last lesson of chapter 436: I can say a laboratory, an experiment, a survey, a discovery and scientific.
\end{chapteropening}

\section[\texorpdfstring{\ml{ലബോറട്ടറി} (labōṟaṭṭaṟi)}{ലബോറട്ടറി (labōṟaṭṭaṟi)}]{\ml{ലബോറട്ടറി} (labōṟaṭṭaṟi) — a laboratory}

\label{lesson:ML-C436-laborattari}



\begin{quote}
Before the new one: say the Malayalam for higher education, then the Malayalam for an evaluation.
\end{quote}

\subsection*{You'll want to know: \ml{ലബോറട്ടറി}}
\textbf{\ml{ലബോറട്ടറി}} — \emph{labōṟaṭṭaṟi} — \textquotedblleft{}a laboratory\textquotedblright{}.

\begin{quote}
\ml{കുട്ടികൾ} \ml{ഉച്ചയ്ക്ക്} \ml{ശേഷം} \ml{ലബോറട്ടറിയിൽ} \ml{പോയി}.
\end{quote}

\begin{quote}
\emph{kuṭṭikaḷ uccaykkŭ śēṣaṁ labōṟaṭṭaṟiyil pōyi.}
\end{quote}

\begin{quote}
The children went to the laboratory in the afternoon.
\end{quote}

In talk often shortened to \ml{ലാബ്}.

\begin{cousinweb}
\textbf{\ml{ലബോറട്ടറി}} (\emph{labōṟaṭṭaṟi}) was borrowed from English \emph{laboratory}.
\end{cousinweb}

\begin{grammarlens}[title={What you've built}]
One more word for giving reasons and opinions.
\end{grammarlens}

\subsection*{Your turn}
\begin{itemize}
\raggedright
  \item \emph{Say it:} \emph{labōṟaṭṭaṟi}
  \item \emph{Say it:} \emph{labōṟaṭṭaṟi}, once more
  \item \emph{Say it:} \emph{kuṭṭikaḷ uccaykkŭ śēṣaṁ labōṟaṭṭaṟiyil pōyi.} — The children went to the laboratory in the afternoon.
  \item \emph{Recall:} say the Malayalam for higher education, then the Malayalam for an evaluation, then say \emph{labōṟaṭṭaṟi} again
\end{itemize}

\begin{tcolorbox}[breakable,colback=teal!4,colframe=teal!35!black,title={Before you move on}]
What does \ml{ലബോറട്ടറി} mean? (\textquotedblleft{}A laboratory\textquotedblright{}.) Say it once more.
\end{tcolorbox}
\section[\texorpdfstring{\ml{പരീക്ഷണം} (parīkṣaṇaṁ)}{പരീക്ഷണം (parīkṣaṇaṁ)}]{\ml{പരീക്ഷണം} (parīkṣaṇaṁ) — an experiment}

\label{lesson:ML-C436-pariksanam}



\begin{quote}
Before the new one: say the Malayalam for an evaluation, then the Malayalam for a laboratory.
\end{quote}

\subsection*{You'll want to know: \ml{പരീക്ഷണം}}
\textbf{\ml{പരീക്ഷണം}} — \emph{parīkṣaṇaṁ} — \textquotedblleft{}an experiment\textquotedblright{}.

\begin{quote}
\ml{ആദ്യത്തെ} \ml{പരീക്ഷണം} \ml{വിജയിച്ചില്ല}.
\end{quote}

\begin{quote}
\emph{ādyatte parīkṣaṇaṁ vijayiccilla.}
\end{quote}

\begin{quote}
The first experiment did not succeed.
\end{quote}

To experiment is \ml{പരീക്ഷണം} \ml{നടത്തുക}.

\begin{cousinweb}
\textbf{\ml{പരീക്ഷണം}} (\emph{parīkṣaṇaṁ}) is borrowed from Sanskrit \emph{parīkṣaṇa}, \textquotedblleft{}trying, testing, an experiment\textquotedblright{} (Monier-Williams), a sister of \ml{പരീക്ഷ}, an exam.
\end{cousinweb}

\begin{grammarlens}[title={What you've built}]
One more word for giving reasons and opinions.
\end{grammarlens}

\subsection*{Your turn}
\begin{itemize}
\raggedright
  \item \emph{Say it:} \emph{parīkṣaṇaṁ}
  \item \emph{Say it:} \emph{parīkṣaṇaṁ}, once more
  \item \emph{Say it:} \emph{ādyatte parīkṣaṇaṁ vijayiccilla.} — The first experiment did not succeed.
  \item \emph{Recall:} say the Malayalam for an evaluation, then the Malayalam for a laboratory, then say \emph{parīkṣaṇaṁ} again
\end{itemize}

\begin{tcolorbox}[breakable,colback=teal!4,colframe=teal!35!black,title={Before you move on}]
What does \ml{പരീക്ഷണം} mean? (\textquotedblleft{}An experiment\textquotedblright{}.) Say it once more.
\end{tcolorbox}
\section[\texorpdfstring{\ml{സർവേ} (sarvē)}{സർവേ (sarvē)}]{\ml{സർവേ} (sarvē) — a survey}

\label{lesson:ML-C436-sarve}



\begin{quote}
Before the new one: say the Malayalam for a laboratory, then the Malayalam for an experiment.
\end{quote}

\subsection*{You'll want to know: \ml{സർവേ}}
\textbf{\ml{സർവേ}} — \emph{sarvē} — \textquotedblleft{}a survey\textquotedblright{}.

\begin{quote}
\ml{വിദ്യാർത്ഥികൾ} \ml{ഗ്രാമത്തിൽ} \ml{ഒരു} \ml{സർവേ} \ml{നടത്തി}.
\end{quote}

\begin{quote}
\emph{vidyārtthikaḷ grāmattil oru sarvē naṭatti.}
\end{quote}

\begin{quote}
The students carried out a survey in the village.
\end{quote}

A survey is \ml{നടത്തുക}, carried out.

\begin{cousinweb}
\textbf{\ml{സർവേ}} (\emph{sarvē}) was borrowed from English \emph{survey}.
\end{cousinweb}

\begin{grammarlens}[title={What you've built}]
One more word for giving reasons and opinions.
\end{grammarlens}

\subsection*{Your turn}
\begin{itemize}
\raggedright
  \item \emph{Say it:} \emph{sarvē}
  \item \emph{Say it:} \emph{sarvē}, once more
  \item \emph{Say it:} \emph{vidyārtthikaḷ grāmattil oru sarvē naṭatti.} — The students carried out a survey in the village.
  \item \emph{Recall:} say the Malayalam for a laboratory, then the Malayalam for an experiment, then say \emph{sarvē} again
\end{itemize}

\begin{tcolorbox}[breakable,colback=teal!4,colframe=teal!35!black,title={Before you move on}]
What does \ml{സർവേ} mean? (\textquotedblleft{}A survey\textquotedblright{}.) Say it once more.
\end{tcolorbox}
\section[\texorpdfstring{\ml{കണ്ടുപിടിത്തം} (kaṇṭupiṭittaṁ)}{കണ്ടുപിടിത്തം (kaṇṭupiṭittaṁ)}]{\ml{കണ്ടുപിടിത്തം} (kaṇṭupiṭittaṁ) — a discovery; an invention}

\label{lesson:ML-C436-kantupitittam}



\begin{quote}
Before the new one: say the Malayalam for an experiment, then the Malayalam for a survey.
\end{quote}

\subsection*{You'll want to know: \ml{കണ്ടുപിടിത്തം}}
\textbf{\ml{കണ്ടുപിടിത്തം}} — \emph{kaṇṭupiṭittaṁ} — \textquotedblleft{}a discovery; an invention\textquotedblright{}.

\begin{quote}
\ml{വൈദ്യുതി} \ml{ഒരു} \ml{വലിയ} \ml{കണ്ടുപിടിത്തമായിരുന്നു}.
\end{quote}

\begin{quote}
\emph{vaidyuti oru valiya kaṇṭupiṭittamāyirunnu.}
\end{quote}

\begin{quote}
Electricity was a great discovery.
\end{quote}

The verb is \ml{കണ്ടുപിടിക്കുക}, to find out, to invent.

\begin{cousinweb}
\textbf{\ml{കണ്ടുപിടിത്തം}} (\emph{kaṇṭupiṭittaṁ}) is a native Malayalam compound: \ml{കണ്ടു}, \textquotedblleft{}having seen\textquotedblright{}, from \ml{കാണുക} (DEDR 1443), and \ml{പിടിത്തം}, \textquotedblleft{}a grasp\textquotedblright{}, from \ml{പിടിക്കുക}, \textquotedblleft{}to catch\textquotedblright{} (DEDR 4148): to see a thing and catch hold of it.
\end{cousinweb}

\begin{grammarlens}[title={What you've built}]
One more word for giving reasons and opinions.
\end{grammarlens}

\subsection*{Your turn}
\begin{itemize}
\raggedright
  \item \emph{Say it:} \emph{kaṇṭupiṭittaṁ}
  \item \emph{Say it:} \emph{kaṇṭupiṭittaṁ}, once more
  \item \emph{Say it:} \emph{vaidyuti oru valiya kaṇṭupiṭittamāyirunnu.} — Electricity was a great discovery.
  \item \emph{Recall:} say the Malayalam for an experiment, then the Malayalam for a survey, then say \emph{kaṇṭupiṭittaṁ} again
\end{itemize}

\begin{tcolorbox}[breakable,colback=teal!4,colframe=teal!35!black,title={Before you move on}]
What does \ml{കണ്ടുപിടിത്തം} mean? (\textquotedblleft{}A discovery; an invention\textquotedblright{}.) Say it once more.
\end{tcolorbox}
\section[\texorpdfstring{\ml{ശാസ്ത്രീയമായ} (śāstrīyamāya)}{ശാസ്ത്രീയമായ (śāstrīyamāya)}]{\ml{ശാസ്ത്രീയമായ} (śāstrīyamāya) — scientific}

\label{lesson:ML-C436-sastriyamaya}



\begin{quote}
Before the new one: say the Malayalam for a survey, then the Malayalam for a discovery.
\end{quote}

\subsection*{You'll want to know: \ml{ശാസ്ത്രീയമായ}}
\textbf{\ml{ശാസ്ത്രീയമായ}} — \emph{śāstrīyamāya} — \textquotedblleft{}scientific\textquotedblright{}.

\begin{quote}
\ml{അവർ} \ml{ശാസ്ത്രീയമായ} \ml{രീതിയിൽ} \ml{കൃഷി} \ml{ചെയ്യുന്നു}.
\end{quote}

\begin{quote}
\emph{avar śāstrīyamāya rītiyil kṛṣi ceyyunnu.}
\end{quote}

\begin{quote}
They farm in a scientific way.
\end{quote}

The adverb is \ml{ശാസ്ത്രീയമായി}, scientifically.

\begin{cousinweb}
\textbf{\ml{ശാസ്ത്രീയമായ}} (\emph{śāstrīyamāya}) is built from Sanskrit \emph{śāstrīya}, \textquotedblleft{}taught in the treatises\textquotedblright{} (Monier-Williams), with the adjective ending -\ml{ആയ}.
\end{cousinweb}

\begin{grammarlens}[title={What you've built}]
One more word for giving reasons and opinions.
\end{grammarlens}

\subsection*{Your turn}
\begin{itemize}
\raggedright
  \item \emph{Say it:} \emph{śāstrīyamāya}
  \item \emph{Say it:} \emph{śāstrīyamāya}, once more
  \item \emph{Say it:} \emph{avar śāstrīyamāya rītiyil kṛṣi ceyyunnu.} — They farm in a scientific way.
  \item \emph{Recall:} say the Malayalam for a survey, then the Malayalam for a discovery, then say \emph{śāstrīyamāya} again
\end{itemize}

\begin{tcolorbox}[breakable,colback=teal!4,colframe=teal!35!black,title={Before you move on}]
What does \ml{ശാസ്ത്രീയമായ} mean? (\textquotedblleft{}Scientific\textquotedblright{}.) Say it once more.
\end{tcolorbox}
\section[(dialogue)]{Words from chapters 429–436 — a review}

\label{lesson:ML-R436-review-429-436}



\begin{quote}
Say the Malayalam for welfare and a ration.
\end{quote}

\subsection*{Your turn}
Say these aloud:

\begin{itemize}
\raggedright
  \item welfare, a ration, financial aid, a benefit, to distribute
  \item corruption, a bribe, to exploit, transparent, to investigate
  \item development, to develop, to modernise, to make headway, a system
  \item research, a researcher, a thesis, a professor, a lecture
  \item admission, an entrance exam, a course of study, a semester, a syllabus
  \item a curriculum, attendance, a hostel, a scholarship, to pass an exam
  \item a rank, a graduate, a graduation ceremony, higher education, an evaluation
  \item a laboratory, an experiment, a survey, a discovery, scientific
\end{itemize}

\begin{tcolorbox}[breakable,colback=teal!4,colframe=teal!35!black,title={Before you move on}]
Welfare? (\textbf{\ml{ക്ഷേമം}}, \emph{kṣēmaṁ}.) A ration? (\textbf{\ml{റേഷൻ}}, \emph{ṟēṣan}.) Financial aid? (\textbf{\ml{ധനസഹായം}}, \emph{dhanasahāyaṁ}.) A benefit? (\textbf{\ml{ആനുകൂല്യം}}, \emph{ānukūlyaṁ}.) To distribute? (\textbf{\ml{വിതരണം} \ml{ചെയ്യുക}}, \emph{vitaraṇaṁ ceyyuka}.) Corruption? (\textbf{\ml{അഴിമതി}}, \emph{aḻimati}.) A bribe? (\textbf{\ml{കൈക്കൂലി}}, \emph{kaikkūli}.) To exploit? (\textbf{\ml{ചൂഷണം} \ml{ചെയ്യുക}}, \emph{cūṣaṇaṁ ceyyuka}.) Transparent? (\textbf{\ml{സുതാര്യമായ}}, \emph{sutāryamāya}.) To investigate? (\textbf{\ml{അന്വേഷിക്കുക}}, \emph{anvēṣikkuka}.) Development? (\textbf{\ml{വികസനം}}, \emph{vikasanaṁ}.) To develop? (\textbf{\ml{വികസിപ്പിക്കുക}}, \emph{vikasippikkuka}.) To modernise? (\textbf{\ml{നവീകരിക്കുക}}, \emph{navīkarikkuka}.) To make headway? (\textbf{\ml{പുരോഗമിക്കുക}}, \emph{purōgamikkuka}.) A system? (\textbf{\ml{സമ്പ്രദായം}}, \emph{sampradāyaṁ}.) Research? (\textbf{\ml{ഗവേഷണം}}, \emph{gavēṣaṇaṁ}.) A researcher? (\textbf{\ml{ഗവേഷകൻ}}, \emph{gavēṣakan}.) A thesis? (\textbf{\ml{പ്രബന്ധം}}, \emph{prabandhaṁ}.) A professor? (\textbf{\ml{പ്രൊഫസർ}}, \emph{prophasar}.) A lecture? (\textbf{\ml{പ്രഭാഷണം}}, \emph{prabhāṣaṇaṁ}.) Admission? (\textbf{\ml{പ്രവേശനം}}, \emph{pravēśanaṁ}.) An entrance exam? (\textbf{\ml{പ്രവേശനപരീക്ഷ}}, \emph{pravēśanaparīkṣa}.) A course of study? (\textbf{\ml{കോഴ്സ്}}, \emph{kōḻsŭ}.) A semester? (\textbf{\ml{സെമസ്റ്റർ}}, \emph{semasṟṟar}.) A syllabus? (\textbf{\ml{സിലബസ്}}, \emph{silabasŭ}.) A curriculum? (\textbf{\ml{പാഠ്യപദ്ധതി}}, \emph{pāṭhyapaddhati}.) Attendance? (\textbf{\ml{ഹാജർ}}, \emph{hājar}.) A hostel? (\textbf{\ml{ഹോസ്റ്റൽ}}, \emph{hōsṟṟal}.) A scholarship? (\textbf{\ml{സ്കോളർഷിപ്പ്}}, \emph{skōḷarṣippŭ}.) To pass an exam? (\textbf{\ml{പാസാകുക}}, \emph{pāsākuka}.) A rank? (\textbf{\ml{റാങ്ക്}}, \emph{ṟāṅkŭ}.) A graduate? (\textbf{\ml{ബിരുദധാരി}}, \emph{birudadhāri}.) A graduation ceremony? (\textbf{\ml{ബിരുദദാനം}}, \emph{birudadānaṁ}.) Higher education? (\textbf{\ml{ഉന്നതവിദ്യാഭ്യാസം}}, \emph{unnatavidyābhyāsaṁ}.) An evaluation? (\textbf{\ml{മൂല്യനിർണ്ണയം}}, \emph{mūlyanirṇṇayaṁ}.) A laboratory? (\textbf{\ml{ലബോറട്ടറി}}, \emph{labōṟaṭṭaṟi}.) An experiment? (\textbf{\ml{പരീക്ഷണം}}, \emph{parīkṣaṇaṁ}.) A survey? (\textbf{\ml{സർവേ}}, \emph{sarvē}.) A discovery? (\textbf{\ml{കണ്ടുപിടിത്തം}}, \emph{kaṇṭupiṭittaṁ}.) Scientific? (\textbf{\ml{ശാസ്ത്രീയമായ}}, \emph{śāstrīyamāya}.)
\end{tcolorbox}
